% ======================================================================
% Chapter 8. Normalization: canonical trace functionals and stable ledgers
% Annals+++ / DIAMOND rules enforced:
%   - OBJ-LOCK: the trace functional is a fixed object; auxiliary architecture is tracked by a ledger.
%   - LOCKS: Fourier/window/trace-regularization conventions are inherited from Ch. 5–7 and never altered silently.
%   - CANONICALITY: normalization removes only identity-sector ambiguity; everything else is quantified in a finite ledger.
%   - PORTABILITY: later arguments may use only (T_\Gamma, E_\Gamma) + explicit identity correction.
% ======================================================================

\chapter{Normalization: canonical trace functionals and stable error ledgers}
\label{chap:08-normalization}

\epigraph{Normalization is the act of turning a formula into an object.}{}

% ======================================================================
\section{Purpose and scope}\label{sec:08.scope}
% ======================================================================

\noindent\textbf{Objective.}
The localized trace formalism developed in Chapters~5--7 produces a \emph{raw} (regularized) trace functional
\[
\mathcal L_\Gamma[h]\;=\;\Tr^{\mathrm{reg}}\!\big(h(A)\big),
\qquad A=\sqrt{\Delta-\tfrac14},
\]
whose construction depends on admissible auxiliary architecture (cutoffs, microlocal partitions, quantization conventions,
and in the noncompact case the fixed cusp truncation scheme). The purpose of this chapter is to isolate a normalization layer
which makes the trace functional \emph{canonical} and \emph{portable}:

\begin{itemize}
\item a normalized functional
\[
\mathcal T_\Gamma[h],
\]
defined on an admissible test class, and
\item a consolidated \emph{error ledger}
\[
\mathcal E_\Gamma[h],
\]
which bounds the total instability under all admissible auxiliary changes.
\end{itemize}

\medskip
\noindent\textbf{Design principle.}
Normalization must:
(i) remove only the \emph{identity-sector ambiguity} (the unique part that can drift under harmless architectural changes),
(ii) preserve admissibility of tests, and
(iii) reduce all remaining dependence on auxiliary choices to a \emph{finite list of stable error contributions}
recorded in $\mathcal E_\Gamma[h]$.

\medskip
\noindent\textbf{Main stability statement (informal).}
Once normalization is imposed, any admissible modification of the localization architecture changes the trace functional only by
(i) an explicit identity-sector correction and
(ii) a remainder controlled by the global ledger.

% ======================================================================
\section{Admissible test functions and canonical identity subtraction}
\label{sec:08.test-class}
% ======================================================================

\subsection{Admissible tests and even/odd splitting}\label{subsec:08.even-odd}

Let $h:\R\to\C$ be an admissible test function in the sense fixed in Chapters~2--7
(e.g.\ Schwartz class; or Paley--Wiener with the Fourier conventions locked in Chapter~7).
We use the standard decomposition
\[
h^{\mathrm{ev}}(t)=\frac12\bigl(h(t)+h(-t)\bigr),
\qquad
h^{\mathrm{odd}}(t)=\frac12\bigl(h(t)-h(-t)\bigr),
\qquad h=h^{\mathrm{ev}}+h^{\mathrm{odd}}.
\]
In the untwisted Selberg trace formula only the even part contributes to the geometric side, but we keep the full class for
compatibility with microlocal controls and twisted observables introduced later.

\subsection{The identity-sector ambiguity}\label{subsec:08.identity-ambiguity}

The raw functional $\mathcal L_\Gamma[h]$ contains an identity-sector contribution determined by local behavior of $h$ at $0$
(and, in noncompact settings, by the fixed truncation ledger). This is the \emph{only} source of non-canonical dependence that
cannot be eliminated without explicitly choosing a reference subtraction. The normalization below removes precisely this mode
and nothing else.

\subsection{Canonical subtraction map}\label{subsec:08.zero-mode}

\begin{definition}[Admissible cutoff]\label{def:08.adm-cutoff}
A function $\chi\in C_c^\infty(\R)$ is an \emph{admissible cutoff} if $\chi\equiv 1$ in a neighborhood of $0$.
\end{definition}

\begin{definition}[Normalization map]\label{def:08.norm-map}
Fix an admissible cutoff $\chi$. Define the normalization map
\begin{equation}\label{eq:08.N.map}
\mathsf N_\chi:\; h\longmapsto h^\sharp:=h-h(0)\chi.
\end{equation}
\end{definition}

\begin{remark}[Minimality and admissibility]\label{rem:08.norm-minimal}
The map $\mathsf N_\chi$ removes only the constant mode $h(0)$, and modifies $h$ by a compactly supported smooth term.
Hence it preserves the admissibility class required for the localized trace formula and is compatible with the trace-class
regularization established earlier.
\end{remark}

% ======================================================================
\section{Normalized trace functional}\label{sec:08.normalized-trace}
% ======================================================================

\subsection{Definition}\label{subsec:08.def-first}

\begin{definition}[Normalized trace functional]\label{def:08.norm-trace}
For an admissible test $h$ define
\begin{equation}\label{eq:08.T.def}
\mathcal T_\Gamma[h]
:=\mathcal L_\Gamma[h^\sharp]
=\mathcal L_\Gamma\!\big[h-h(0)\chi\big].
\end{equation}
\end{definition}

\begin{remark}[Immediate properties]\label{rem:08.immediate}
\leavevmode
\begin{enumerate}
\item\label{it:08.immediate1} $\mathcal T_\Gamma[\chi]=0$.
\item\label{it:08.immediate2} $\mathcal T_\Gamma$ is linear on the admissible test class.
\item\label{it:08.immediate3} Dependence on $\chi$ is confined to the identity-sector subtraction.
\end{enumerate}
\end{remark}

\subsection{Cutoff change formula}\label{subsec:08.cutoff-change}

\begin{lemma}[Cutoff change formula]\label{lem:08.cutoff-change}
Let $\chi_1,\chi_2$ be admissible cutoffs and $h$ an admissible test function. Define
\[
\mathcal T^{(j)}_\Gamma[h]:=\mathcal L_\Gamma\big[h-h(0)\chi_j\big],\qquad j=1,2.
\]
Then
\begin{equation}\label{eq:08.cutoff-change}
\mathcal T^{(1)}_\Gamma[h]-\mathcal T^{(2)}_\Gamma[h]
=-\,h(0)\,\mathcal L_\Gamma[\chi_1-\chi_2].
\end{equation}
\end{lemma}

\begin{proof}
By linearity of $\mathcal L_\Gamma$,
\[
\mathcal T^{(1)}_\Gamma[h]-\mathcal T^{(2)}_\Gamma[h]
=\mathcal L_\Gamma[h-h(0)\chi_1]-\mathcal L_\Gamma[h-h(0)\chi_2]
=-h(0)\,\mathcal L_\Gamma[\chi_1-\chi_2].
\]
\end{proof}

\begin{theorem}[Canonicality modulo identity-sector correction]\label{thm:08.canonical-mod-identity}
Assume the identity contribution in the localized trace decomposition of Chapters~6--7 is explicit in the fixed convention.
Then for any admissible cutoffs $\chi_1,\chi_2$ and any admissible test $h$,
the difference $\mathcal T^{(1)}_\Gamma[h]-\mathcal T^{(2)}_\Gamma[h]$ is an explicit identity-sector correction depending only on
$h(0)$ and on the fixed identity term of the trace formula.
\end{theorem}

\begin{proof}
By Lemma~\ref{lem:08.cutoff-change}, the difference equals $-h(0)\mathcal L_\Gamma[\chi_1-\chi_2]$.
Since $\chi_1-\chi_2$ is compactly supported, smooth, and vanishes near $0$, its contribution is confined to the identity sector
in the Selberg/wave-trace expansion in the fixed convention; hence the explicit identity term yields an explicit formula.
\end{proof}

% ======================================================================
\section{Normalization stability under admissible auxiliary changes}\label{sec:08.stability}
% ======================================================================

\noindent
The construction of $\mathcal L_\Gamma$ may depend on admissible auxiliary architecture:
quantization convention, microlocal partition, and (in the noncompact case) regularization of the continuous spectrum.
We formalize stability by isolating each dependence into a bounded ledger contribution.

\subsection{Quantization stability ledger}\label{subsec:08.quantization}

Let $\Op$ denote the quantization used to define microlocal projectors and localized kernels (Chapter~5).
Different admissible quantizations differ by lower-order operators; hence their induced trace functionals differ by controlled remainders.

\begin{assumption}[Quantization stability]\label{ass:08.quant-stability}
Let $\Op_1,\Op_2$ be admissible quantizations, and let $\mathcal L^{(1)}_\Gamma,\mathcal L^{(2)}_\Gamma$ be the corresponding raw traces.
Then for every admissible test $g$,
\begin{equation}\label{eq:08.quant.err}
\mathcal L^{(1)}_\Gamma[g]-\mathcal L^{(2)}_\Gamma[g]=\Err_{\mathrm{quant}}(g),
\end{equation}
where $\Err_{\mathrm{quant}}(g)$ satisfies the microlocal remainder bounds proved in Chapter~5.
\end{assumption}

\begin{definition}[Quantization ledger]\label{def:08.quant-ledger}
For admissible $h$ define
\begin{equation}\label{eq:08.E.quant}
\mathcal E^{\mathrm{quant}}_\Gamma[h]
:=\sup_{\Op_1,\Op_2\ \mathrm{admissible}} \big|\Err_{\mathrm{quant}}(h^\sharp)\big|.
\end{equation}
\end{definition}

\subsection{Microlocal refinement stability ledger}\label{subsec:08.refinement}

Let $\{\Pi_j\}$ be an admissible microlocal partition of unity used in the localization procedure (Chapters~4--5).
Refinements change only the decomposition, not the underlying operator, and introduce controlled remainder terms.

\begin{assumption}[Refinement stability]\label{ass:08.refine-stability}
Let $\{\Pi_j\}$ and $\{\widetilde\Pi_k\}$ be admissible partitions, with $\{\widetilde\Pi_k\}$ a refinement of $\{\Pi_j\}$.
Let $\mathcal L_{\Gamma,\{\Pi_j\}}$ and $\mathcal L_{\Gamma,\{\widetilde\Pi_k\}}$ be the corresponding raw traces.
Then for every admissible test $g$,
\begin{equation}\label{eq:08.ref.err}
\mathcal L_{\Gamma,\{\Pi_j\}}[g]-\mathcal L_{\Gamma,\{\widetilde\Pi_k\}}[g]=\Err_{\mathrm{ref}}(g),
\end{equation}
where $\Err_{\mathrm{ref}}(g)$ satisfies the microlocal remainder bounds proved in Chapter~5.
\end{assumption}

\begin{definition}[Refinement ledger]\label{def:08.ref-ledger}
For admissible $h$ define
\begin{equation}\label{eq:08.E.ref}
\mathcal E^{\mathrm{ref}}_\Gamma[h]
:=\sup_{\text{admissible refinements}} \big|\Err_{\mathrm{ref}}(h^\sharp)\big|.
\end{equation}
\end{definition}

\subsection{Tail and continuous-spectrum ledgers}\label{subsec:08.tail-cont}

Localization architectures may involve time/length truncations (tails) and, in noncompact settings, a continuous-spectrum
regularization scheme (Chapter~7). These contributions are bounded by the sharpest estimates established earlier.

\begin{definition}[Tail and continuous-spectrum ledgers]\label{def:08.tail-cont-ledgers}
Let $\mathcal E^{\mathrm{tail}}_\Gamma[h]$ denote the best available bound for tail contributions in the localization procedure applied to $h^\sharp$.
If a continuous spectrum term is present, let $\mathcal E^{\mathrm{cont}}_\Gamma[h]$ denote the corresponding bound for the continuous-spectrum
regularization remainder in the fixed convention.
Both quantities are taken from the estimates proved in Chapters~6--7.
\end{definition}

% ======================================================================
\section{Consolidated normalization ledger}\label{sec:08.consolidated-ledger}
% ======================================================================

\begin{definition}[Total normalization ledger]\label{def:08.total-ledger}
For admissible $h$ define the consolidated normalization ledger
\begin{equation}\label{eq:08.E.total}
\mathcal E_\Gamma[h]
:=\mathcal E^{\mathrm{quant}}_\Gamma[h]
+\mathcal E^{\mathrm{ref}}_\Gamma[h]
+\mathcal E^{\mathrm{tail}}_\Gamma[h]
+\mathcal E^{\mathrm{cont}}_\Gamma[h].
\end{equation}
\end{definition}

\begin{remark}[Finiteness and portability]\label{rem:08.ledger.portable}
The point of \eqref{eq:08.E.total} is not to optimize constants, but to compress all admissible architectural dependence into a
finite list of stable terms. Later arguments may therefore treat $\mathcal E_\Gamma[h]$ as a black-box bound with known provenance.
\end{remark}

% ======================================================================
\section{Normalization stability theorem}\label{sec:08.norm-stability}
% ======================================================================

\begin{theorem}[Stability of the normalized trace]\label{thm:08.normalization-stability}
Fix an admissible cutoff $\chi$ and define $\mathcal T_\Gamma[h]=\mathcal L_\Gamma[h-h(0)\chi]$.
Let $\mathcal T'_\Gamma[h]$ be the normalized trace obtained from any other admissible construction
(e.g.\ a different admissible cutoff $\chi'$, different admissible quantization, microlocal partition,
or continuous-spectrum regularization), with corresponding consolidated ledgers $\mathcal E_\Gamma$ and $\mathcal E'_\Gamma$.
Then for every admissible test $h$,
\begin{equation}\label{eq:08.stability.main}
\mathcal T_\Gamma[h]-\mathcal T'_\Gamma[h]
=
\Delta_{\mathrm{id}}(h;\chi,\chi')
+\Err_{\mathrm{norm}}(h),
\end{equation}
where:
\begin{enumerate}
\item $\Delta_{\mathrm{id}}(h;\chi,\chi')$ is an explicit identity-sector correction depending only on $h(0)$ and the fixed identity term;
\item the remainder satisfies
\begin{equation}\label{eq:08.stability.err}
|\Err_{\mathrm{norm}}(h)|
\le \mathcal E_\Gamma[h]+\mathcal E'_\Gamma[h].
\end{equation}
\end{enumerate}
In particular, modulo the explicit identity correction, the normalized trace functional is canonical up to a quantified error
controlled by the total ledgers.
\end{theorem}

\begin{proof}
The cutoff dependence is governed by Lemma~\ref{lem:08.cutoff-change}, yielding $\Delta_{\mathrm{id}}$ via
Theorem~\ref{thm:08.canonical-mod-identity}.
All remaining admissible architectural changes are controlled by Assumptions~\ref{ass:08.quant-stability} and~\ref{ass:08.refine-stability}
and by the tail/continuous-spectrum bounds recorded in Definition~\ref{def:08.tail-cont-ledgers}.
Summing these bounds gives \eqref{eq:08.stability.err}.
\end{proof}

% ======================================================================
\section{Consequences for later chapters}\label{sec:08.consequences}
% ======================================================================

\begin{corollary}[Drop-in principle]\label{cor:08.drop-in}
All arguments in Chapters~9--10 which depend on the localized trace formula may be formulated solely in terms of
the normalized functional $\mathcal T_\Gamma[h]$ and the consolidated ledger $\mathcal E_\Gamma[h]$,
without explicit reference to the auxiliary architecture used to define the raw trace, provided the identity-sector
correction $\Delta_{\mathrm{id}}$ is treated explicitly.
\end{corollary}

\begin{corollary}[Stability under rescaling of tests]\label{cor:08.rescaling}
Let $h_T(t):=h(t/T)$ with $T\ge 1$, and assume $h_T$ remains admissible for the range of $T$ under consideration.
Then the map
\[
T\longmapsto \mathcal T_\Gamma[h_T]
\]
is stable under admissible auxiliary changes, with stability constants controlled by $\mathcal E_\Gamma[h_T]$ uniformly in $T$
(modulo the explicit identity-sector correction).
\end{corollary}

% ======================================================================
\section{Bibliographic remarks}\label{sec:08.bib}
% ======================================================================

Normalization by explicit subtraction of the identity-sector ambiguity is standard in microlocal trace theory and in the Selberg
trace formula tradition. For background on spectral geometry and the Selberg trace formula see
\cite{Buser1992,Hejhal1983,Iwaniec2002}. For semiclassical trace methods and scattering resonance technology see
\cite{Zworski2012,DyatlovZworski2019}.

% ======================================================================
% Audit (Chapter 8, normalization layer)
% ======================================================================

\section*{Audit (Normalization layer)}\label{sec:08.audit}

\noindent\textbf{Locks.}
\begin{itemize}
\item[(L8.1)] OBJ-LOCK: $\mathcal L_\Gamma[h]=\Tr^{\mathrm{reg}}(h(A))$ with the fixed $A=\sqrt{\Delta-\tfrac14}$.
\item[(L8.2)] Normalization map: $h^\sharp=h-h(0)\chi$ with $\chi\equiv1$ near $0$.
\item[(L8.3)] Cutoff dependence is explicit: \eqref{eq:08.cutoff-change}.
\item[(L8.4)] All auxiliary dependence beyond identity-sector is ledgered in $\mathcal E_\Gamma[h]$.
\end{itemize}

\medskip
\noindent\textbf{Failure modes.}
\begin{itemize}
\item[(F8.1)] Subtracting a nonlocal term (not compactly supported) $\Rightarrow$ changes hyperbolic/orbit content (invalid normalization).
\item[(F8.2)] Mixing cusp conventions in $\Tr^{\mathrm{reg}}$ $\Rightarrow$ ledger mismatch with Chapter~7.
\item[(F8.3)] Treating ledger bounds as ``zero'' $\Rightarrow$ portability claims become false.
\end{itemize}

\medskip
\noindent\textbf{STATUS\_CERT.}
CLOSED (structural): normalization functional $\mathcal T_\Gamma$ and consolidated ledger $\mathcal E_\Gamma$ are defined and
their dependence on admissible architecture is reduced to an explicit identity correction plus quantified ledger bounds.
NEXT(1): instantiate $\mathcal E^{\mathrm{tail}}_\Gamma$ and $\mathcal E^{\mathrm{cont}}_\Gamma$ with the explicit bounds proved in Chapters~6--7
for the specific window family $h_{\lambda,\eta}$, and record the resulting uniform $(\lambda,\eta)$-scaling of $\mathcal E_\Gamma[h_{\lambda,\eta}]$.

% ======================================================================
% End of 08-normalization.tex
% ======================================================================
